% Part III-C, discussion (included by main_en.tex).
% TODO when rlvr seed 3 is in: check the paragraph on the tree shape against its accuracies.

\section{Discussion}\label{sec:discussion}

\paragraph{Within a problem, the judge sees little of where a solution is going.} Offline, \kev{}'s \emph{on track}
score separates the correct from the incorrect continuations of one problem's 128-token prefixes with an AUROC of
\NUM{aucw_128_kevtrack}, no better than the AR judge of the same size (\NUM{aucw_128_artrack}) or the policy's own
log-probability (\NUM{aucw_128_logprob}), which costs nothing. Across problems every signal does better, because each
partly sees how hard the problem is, and that does not help to choose among one problem's prefixes. In training, on
the policy's own roots, the separation was weaker still (\NUM{auditotaucw_all_treekev}
\ci{\NUM{auditotaucwlo_all_treekev}}{\NUM{auditotaucwhi_all_treekev}}), and keeping \kev{}'s two roots of four raised
the share of correct leaves by \NUM{auditgain_all_treekev}
\ci{\NUM{auditgainlo_all_treekev}}{\NUM{auditgainhi_all_treekev}} points, against \NUM{auditgain_all_treerand} for a
random choice and \NUM{auditbest_all_treekev} for the best two in hindsight. The gain was larger early in training
(\NUM{auditgain_firsthalf_treekev} points in the first half, \NUM{auditgain_secondhalf_treekev} in the second, 20
audited steps each, without intervals), while the policy's roots came to lead to correct answers more often and the
room for a choice shrank (best two in hindsight: \NUM{auditbest_firsthalf_treekev} and
\NUM{auditbest_secondhalf_treekev} points).

\paragraph{A small gain at the roots did not change what the policy learned.} With the same tree shape, keeping
\kev{}'s roots instead of random ones left the final accuracy unchanged within the intervals over problems on four of
the five benchmarks; on GSM8K it was \NUM{dvtreerand_treekev_gsm}
\ci{\NUM{dvtreerandlo_treekev_gsm}}{\NUM{dvtreerandhi_treekev_gsm}} points higher, one interval of five over two seeds,
which we do not count as an effect. The idea behind the judge's choice is that it gives GRPO more groups with mixed
outcomes, the only groups with a learning signal. It did not: in the last 20 steps \NUM{allwrong_treekev} of the
groups had no correct leaf under \texttt{tree\_kev} and \NUM{allwrong_treerand} under \texttt{tree\_random}, and
\NUM{allright_treekev} and \NUM{allright_treerand} had only correct ones. The choice has a cost in principle: the kept
roots are no longer samples of the policy, so the gradient is no longer an unbiased estimate of the policy gradient.
We saw no effect of that on accuracy. One seed of \texttt{tree\_kev} came to repeat its solution after the first
answer, which neither seed of \texttt{tree\_random} did. Its reward is the verifier's alone, which does not penalise a
repetition that restates the answer, so this is not the judge paying for repetition; whether the choice of roots
caused it, one seed of two cannot tell.

\paragraph{Whether the tree shape helps is open.} Both tree conditions end above flat sampling in the seed mean, by
\NUM{d_treerand_math} and \NUM{d_treekev_math} points on MATH-500 and \NUM{d_treerand_aime24} and
\NUM{d_treekev_aime24} on AIME24, and \texttt{tree\_random} generated fewer tokens per step than flat sampling
(\NUM{gentok_treerand} thousand against \NUM{gentok_rlvr} thousand). But each tree condition has one seed well above
\texttt{rlvr} and one at its level, and at step 100 neither was ahead. The intervals over problems do not include the
variation between seeds, and with two seeds per condition the difference may be the seeds'. Testing the tree shape
needs more seeds than we ran.

\paragraph{Limitations.} One policy size (4B), one judge question (\emph{on track}) asked of 128-token roots, one tree
shape (four roots, two kept, four leaves each), 200 steps at the learning rate of Part II, and \NUM{seeds_treekev}
seeds per tree condition (\texttt{rlvr}: \NUM{seeds_rlvr}). With more roots per prompt a judge has more to choose from,
and so do its errors. The audit continues all roots only every tenth step, 20 audited steps per run; its intervals
resample those steps, not seeds. The offline study used the base policy's prefixes, not the training policy's. We
report tokens, not time: the runs shared GPUs with other experiments.

\paragraph{Relation to the earlier versions.} Versions 1--3 of this paper reported pruning speed-ups, KV-cache reuse
rates, per-node latencies and theorem-proving results (MiniF2F-Lean4, ProofNet) that were not measured (see the notice
at the start and \file{CORRECTION.md}). None of these is measured here: we count tokens, not time, and ran no theorem
proving. What we could test, we did. The pruning rule of v3, prune a prefix when the judge's $p(\text{yes})<0.3$,
prunes \NUM{thrpruned_128_kevtrack} of the 128-token prefixes with the \emph{on track} question and catches
\NUM{thrcaught_128_kevtrack} of the dead ones (\S\ref{sec:res-prefix}). Choosing which roots to continue by \kev{}, the
idea behind those versions, did not make GRPO learn better than a random choice of roots.
